% F9 -- The hour token: how it enters, and the ablation that tests whether it is read.
\begin{figure}[t]
\centering
% ---- (a) the 24-hour weight map ----------------------------------------
\fitwidth{%
\begin{tikzpicture}[font=\scriptsize,
  cell/.style={draw=clrule,fill=clsoft,minimum width=0.55cm,minimum height=0.50cm,inner sep=0pt},
  w2/.style={cell,draw=clgain,fill=clgain!25},
  w4/.style={cell,draw=clgain,fill=clgain!55},
]
\foreach \h in {0,1,2,3,6,7,8,9,10,11,12,13,14,15,16,17,20,21,22,23}
  {\node[cell] at (\h*0.58,0) {};}
\foreach \h in {4,5}   {\node[w2] at (\h*0.58,0) {$\times2$};}
\foreach \h in {18,19} {\node[w4] at (\h*0.58,0) {$\times4$};}
\foreach \h/\lbl in {0/00h,4/04h,6/06h,12/12h,18/18h,20/20h,23/23h}
  {\node[anchor=north,text=clrule] at (\h*0.58,-0.32) {\lbl};}
\node[anchor=west,font=\scriptsize,text=clrule,text width=13cm] at (0,0.72)
  {\textbf{(a)} loss weight applied to the target at each hour of the day. Bounds are
   half-open, so ``04--06h'' means hours 4 and 5.};
\end{tikzpicture}}

\vspace{0.45cm}
\replegend{\swatch{clmodel}~real hours \qquad \swatch{clnonsig}~hours shifted $+12$h}
\fitwidth{%
\begin{tikzpicture}
\begin{groupplot}[
  group style={group size=2 by 1, horizontal sep=1.6cm, ylabels at=edge left},
  repbar, width=0.43\linewidth, height=4.9cm,
  ylabel={top-1 accuracy (\%)}, xtick=data,
  every node near coord/.append style={font=\scriptsize},
]
% ---- (b) ablation ----
\nextgroupplot[title={\footnotesize (b) corrupting only the hour token},
  symbolic x coords={ha,hb,hc,hd},
  xticklabels={real,{shuffled\\\textcolor{clloss}{$-3.0\pt$}},{fixed 12h\\\textcolor{clloss}{$-0.7\pt$}},{shift $+12$h\\\textcolor{clgain}{$+0.5\pt$}}},
  x tick label style={font=\scriptsize,align=center},
  ymin=37, ymax=49.5, ytick={38,40,...,48}, /pgf/bar width=11pt]
\addplot[fill=clmodel,draw=none] coordinates {(ha,44.7)(hb,41.7)(hc,44.0)(hd,45.2)};
%
% ---- (c) where the shift lands ----
\nextgroupplot[title={\footnotesize (c) where the $+12$h shift lands},
  symbolic x coords={win,out},
  xticklabels={inside window,outside window},
  x tick label style={font=\scriptsize},
  ymin=36, ymax=51.5, ytick={36,38,...,50}, /pgf/bar width=13pt]
\addplot[fill=clmodel,draw=none] coordinates {(win,42.7)(out,45.3)};
\addplot[fill=clnonsig,draw=none] coordinates {(win,39.7)(out,47.1)};
\end{groupplot}
\end{tikzpicture}}
\caption{The hour-of-day ablation}
\label{fig:hour-ablation}
\figtext{%
Does the model actually read the hour? (a)~Each event enters the sequence as a
dedicated hour token followed by its \cid{cell\_id}, and the number in a
shaded box is the multiplier applied to the target's loss at that hour,
$\times2$ at 04--05h and $\times4$ at 18--19h. Bounds are half-open, so ``04--06h'' means hours 4 and 5, and multiplying the loss duplicates no data: it
only tells the optimiser that a mistake there costs more. (b)~The same trained
checkpoint on the same test users, with the cell sequence untouched and only
the hour token corrupted. The argument is the difference between two of the
bars: replacing every hour by a constant costs $0.7\pt$, while shuffling the
hours, so that the values are still real hours but in the wrong order, costs
$3.0\pt$. If the token merely occupied a slot, destroying its \emph{ordering}
could not cost more than removing it altogether; it costs four times more, so
the model has learned which hour goes with which cell. (c)~Shifting every hour
by twelve moves the aggregate by only $+0.5\pt$, and that near-zero hides a
swap: accuracy inside the transition windows falls $42.7 \rightarrow 39.7\,\%$
while accuracy outside them rises $45.3 \rightarrow 47.1\,\%$. The two halves
of the day exchanged their difficulty along with their labels. \smob{}, 80\,\%
context, $n = 1{,}747$ predictions.}
\end{figure}
